% !TeX root = ../papers/nc-p3.tex
\section{反対位置の階数と二次関係の零点}
\label{sec:opposite-evaluation}
四窓の標準テンソルを反対位置の切断で最小分解し、
$w=\sum_{h=1}^{\rho} A_h\otimes B_h$と書く。
ここで$A_h\in V_0\otimes V_2$, $B_h\in V_1\otimes V_3$であり、
それぞれの組は一次独立である。各位置の非零共変ベクトル$x$で、
対応する$A_h$または$B_h$を評価した$\rho$本の列の階数を$r_i(x)$とする。
隣接対の二次関係$R_i$の共通零点スキームを
$\Gamma_i\subset\mathbb P(V_i^*)\times\mathbb P(V_{i+1}^*)$とする。

\begin{theorem}[AS対象の第三階数は4以上]\label{thm:opposite-rank-at-least-four}
標準AS対象では全ての位置で$r_i(x)\ge2$である。
また各$\Gamma_i$は非空であり、その全ての点$(x,y)$で
$r_i(x)+r_{i+1}(y)\le\rho$である。従って$\rho\ge4$である。
\end{theorem}
\begin{proof}
位置$i$を$x$で縮約した三テンソルは一因子階数4から非零であり、
その中央因子（位置$i+2$）の支持の次元は、反対側の行列の一次独立性から$r_i(x)$である。
系\ref{cor:opposite-singular-rank-two}の中央支持の障害により階数1は不可能である。
一方、$\Gamma_i$はSegre像$\mathbb P^3\times\mathbb P^3\subset\mathbb P^{15}$と
$\mathbb P(R_i^\perp)=\mathbb P^9$の交叉なので、射影次元定理から非空である。
例えば位置0,1を縮約した二テンソルは、列を$A_h(x),B_h(y)$とする
$4\times\rho$行列$P(x),Q(y)$を用いて$P(x)Q(y)^{\mathsf T}$となる。
これが零であることは$R_0$の全元が$(x,y)$で消えることと同値であり、
$\operatorname{im} Q(y)^{\mathsf T}\subset\ker P(x)$から不等式が従う。
他の辺でも同様である。
\end{proof}

\begin{theorem}[第三階数4では全ての隣接点対応が正次元]\label{thm:rho-four-positive-incidence}
標準AS対象で$\rho=4$なら、各$\Gamma_i$の次元は1以上である。
その全ての点では両評価階数が2で、上の像と核は同じ平面になる。
\end{theorem}
\begin{proof}
$\Gamma_0$が有限と仮定する。六本の双一次式の零点は滑らかな六次元の積空間で
余次元6なのでKoszul分解を持つ（零点の外では局所的に単元を含む）。
その分解を$\mathcal O(3,0)$でひねると、外次数$j=1,\ldots,6$の
$\mathcal O(3-j,-j)$は全コホモロジーが零である。
$j=1,2,3$は第二因子、$j=4,5,6$は第一因子の次数が$-1,-2,-3$のいずれかになる。
よって三次式の$\Gamma_0$への制限は単射（実際には同型）である。
非被約な零点にも同じ結論が成立する。

ところが$P(x)Q(y)^{\mathsf T}=0$は$\Gamma_0$上のスキームとしての方程式である。
各点の近傍で$Q(y)$の二次小行列式の一つが単元なので、その転置の像は
階数2の局所直和因子を含む。従って$P(x)$の全三次小行列式が$\Gamma_0$上で消える。
単射性によりこれらは$x$の三次式として恒等的に零である。
従って$P(x)$は$\mathbb P^3$の全点でちょうど階数2となる。

この最後の場合を除く。行列束
$4\mathcal O(-1)\to4\mathcal O$の像を階数2のベクトル束$L$、余核を$N$とする。
超平面類を$h$として$c_1(L)=-a h$, $c_2(L)=b h^2$と書く。
$L^*$と$L(1)$は大域生成なので行列式の次数から$0\le a\le2$である。
$c(N)=1/c(L)$と$\operatorname{rk}N=2$より$a^3-2ab=0$であり、整数$b$に対し$a=1$はない。
$a=0$なら$L^*$、$a=2$なら$L(1)$が行列式自明な大域生成束であり、自明束となる
（一点で基底を与える二切断の外積は非零の定数）。
前者では$L\to4\mathcal O$の像が固定平面となって位置2の一因子階数4に反する。
後者では$4\mathcal O(-1)\to L$が定数行列で二次元の共通核を持ち、
四つの$A_h$の一次独立性に反する。従って$\Gamma_0$は有限ではない。
巡回して全辺について従い、評価の等号は前定理による。
\end{proof}

\begin{corollary}[第三階数5、有限な点対応の三次式制約]\label{cor:rho-five-minor-budget}
$\rho=5$とし、位置$i$の$4\times5$評価行列の三次小行列式が張る三次式空間の次元を$c_i$とする。
標準AS対象で$\Gamma_i$が有限なら$c_i+c_{i+1}\le20$である。
\end{corollary}
\begin{proof}
前定理のKoszul分解により、両因子の三次式はそれぞれ$\Gamma_i$上に単射に制限される。
また有限完全交叉$\Gamma_i$の長さは$(h_1+h_2)^6$の$h_1^3 h_2^3$係数20である。
局所環上で$P: \mathcal O^5\to\mathcal O^4$の二次小行列式の単元を使い、
階数2の恒等ブロックを分裂させる。$PQ^{\mathsf T}=0$より$Q^{\mathsf T}$は
残る三次元の中間空間に入る。さらに$Q^{\mathsf T}$の単元小行列式でその中に
階数2の恒等ブロックを分裂させる。対応する二列では$P$が零である。
残る中間空間は一次元なので、$P,Q^{\mathsf T}$の残りは二行一列$u$と一行二列$v$で、
$uv=0$である。この表示で三次小行列式イデアルはそれぞれ$u$と$v$の成分で生成される。
従って二イデアルの積は零である。

有限スキーム上の直線束を自明化し、二つのイデアルを$I,J$とする。
$H^0(\mathcal O_{\Gamma_i})$はArtin局所完全交叉環の積でFrobeniusである。
従って積とtraceの完全対合で$I,J$は直交し、$\dim I+\dim J\le20$となる。
制限の単射性から$\dim I\ge c_i$, $\dim J\ge c_{i+1}$である。
これは非被約な点も長さに含む。
\end{proof}

\begin{proposition}[階数4の平面の照合と定数ファイバーの全型]\label{prop:rho-four-plane-fibers}
$\rho=4$の標準AS対象を仮定する。位置0,1で係数空間を$H=k^4$とし、
$P:H\to V_2$, $Q^{\mathsf T}:V_3^*\to H$と読む。
評価階数2以下の小行列式スキームを$D_0,D_1$とすると
\[
 \Gamma_0=D_0\mathbin{\times}_{\mathrm{Gr}(2,H)}D_1,qquad
 x\longmapsto\ker P(x),\quad y\longmapsto\operatorname{im}Q(y)^{\mathsf T}.
\]
平面$K\subset H$におけるファイバーは二つの射影線型空間の積である。
片側の射影次元を$a-1$とすると、非空なら$a=1,2,3$である。
この部分で残る二つの係数行列を制限した$4\times a$束は、左右定数基底変更で
$L_e^{\mathsf T}$の直和となり、その添字は順序を除き
\[
 a=1:\ (1,0,0),\qquad a=2:\ (2,0)\ \text{または}\ (1,1),\qquad a=3:\ (3)
\]
で尽くされる。$L_0^{\mathsf T}$は入力が零次元の一行を表す。
\end{proposition}
\begin{proof}
非空な点では両行列の階数が2である。二次の単元小行列式を選べば、
積が零のスキーム上で双方の三次小行列式は全て消え、階数2の像と核が一致する。
逆も局所直和因子の包含から従うのでファイバー積はスキームとして成立する。
$K$を固定すると$P(x)|_K=0$と$Q(y)|_{K^\perp}=0$は各変数について線型である。
例えば第一の解空間を$X$とする。$H/K$の基底に対応する二行列を$A,B:X\to V_2$とすれば、
$0\ne x\in X$で$A x,B x$は独立である。従って全ての$(s:t)$で$sA+tB$は単射である。
Kronecker標準形\cite[Theorem 1]{DDLpencils}では右特異ブロックは常に核を持ち、
非空の正方正則ブロックはある射影点で行列式が消える。
よって左特異ブロックしかなく、その個数は$4-a$、添字の和は$a$となる。
この分割を列挙すると表示の全型を得る。他方も同じである。
\end{proof}

特に次の二つの幾何的可能性を調べればよい。正次元の既約成分上で平面が一定なら、
少なくとも一方で$a\ge2$となり、その側で消える係数の二次元束は
全ての行列が階数2以下の直線をなす。平面が動くなら、係数行列式の四次曲面
$\det(\sum z_h A_h)=0$は動く直線$\mathbb P(K)$を含み、もう一方の行列式も対応する
$\mathbb P(K^\perp)$を含む。恒等的に消える行列式もこの条件に含める。
従って両方が滑らかな四次曲面なら標準AS性は不可能である。
実際、滑らかな行列式曲面上の行列は階数3であり、またその上の直線は
自己交叉$-2$で法束に切断がなく、射影的な直線スキームは有限である。

\subsection{定数平面の排除}
\begin{lemma}[中央二次元の三テンソルの束]\label{lem:middle-binary-pencil}
標準AS対象の三次交叉の元で中央の支持が二次元なら、その二つの端の支持の上で、
中央を縮約する行列束の非零パラメータの全てで階数は端の二次元数の小さい方に等しい。
両端の次元は異なる。
\end{lemma}
\begin{proof}
中央の共変ベクトル$\lambda\ne0$で得た行列が両端に非零の核を持つと、
一方の核で元のテンソルを縮約して中央が$\ker\lambda$に支持される二次関係を得る。
端の空間は実際の支持なのでこの縮約は非零である。反対の核でも同じ中央直線の純粋関係が生じ、
連続純粋関係の禁止に反する。正方なら斉次の行列式に射影零点があるので同様に矛盾する。
\end{proof}

\begin{theorem}[各照合平面での零点対の一意性]\label{thm:matched-plane-unique-pair}
第三階数4の標準AS対象では、実現される各平面$K$の上の
命題\ref{prop:rho-four-plane-fibers}のファイバーは一点である。
対応する二つの係数行列束は、それぞれ零の行または列を除くと
$L_3$または$L_3^{\mathsf T}$である。従って各辺の正次元の点対応では平面そのものが動く。
\end{theorem}
\begin{proof}
位置0の非零解$x$で縮約すると、中央支持は二次元であり、
両端の行列束は$B$の係数平面$K^\perp$そのものである。
共通零行・零列を除き、必要なら転置すると、前補題とKronecker標準形により
この束は$L_e^{\mathsf T}$の直和で、今度は零添字もない。
二つの係数行列が独立で両端は高々四次元なので、その添字の組は
$(1),(2),(1,1),(3)$で尽くされる。位置1の解から$A$の係数平面$K$にも同じ結論を得る。

位置0の解空間$X$が次元$a\ge2$と仮定する。ここで
$K$の行列は$X$を消し、残り二行列の$X$への制限を$C_s,C_t$とする。
命題\ref{prop:rho-four-plane-fibers}により、その非零行のブロックは
$(3),(2),(1,1)$のいずれかである。平面$K^\perp$の束を$D_s,D_t$と書くと、
位置0を$X$で縮約した二つの隣接切断の行列は、零行・零列を除き
$C_s\otimes D_s+C_t\otimes D_t$と、$D_s,D_t$を転置した同じ式である。
両束の標準化はパラメータを固定する厳密同値で行える。
行は$C$、列は$D$のブロック添字として、この二つの階数は順に
\[
\begin{array}{c|rrrr}
 & (1)&(2)&(1,1)&(3)\\\hline
 (3)&3,4&6,8&6,8&9,12\\
 (2)&2,3&4,6&4,6&6,8\\
 (1,1)&2,4&4,6&4,8&6,8
\end{array}
\]
である。これは各$L_e^{\mathsf T}$の二係数を$(e+1)\times e$単位埋込みの
上詰め、下詰めとして代入するだけで再計算できる
（\path{computations/verify_matched_pencils.py}）。
元の切断の階数は6であり、制限の出力のいずれか一因子の支持が真部分空間なら
制限階数は5以下である。等号6なら制限の像が元の像全体となり、一因子階数4に反する。
表の$(1)$列以外では8以上があるか、第二の階数が6でその$D$側の出力が二次元となる。
従って$D$は$(1)$に限られる。

先に$a=3$を除く。このとき$K$の二行列は$X$を消すので位置0に共通の一次元支持を持ち、
$(1)$型の$D$の二行列も位置1または3に共通の一次元支持を持つ。
従って隣接する二位置に固定文字$u,v$があり、テンソルの各項はそのどちらかを持つ。
順序をこの二位置から始め、必要なら逆順として、
その辺の関係$R$は$ku\otimes V+V'\otimes kv$に入る。
$\dim R=6$なので$V'\otimes kv$との交叉は少なくとも$6-3=3$次元である。
最初の文字を$u$の零化空間で縮約すると、隣の位置は$kv$に支持される。
一因子階数4なので非零縮約があり、次の辺には$v\otimes z\ne0$が入る。
従って連続純粋関係となる。これでどの辺のどちら側でも$a=3$は不可能である。
$a=2$でも$(1)$型の$D$は一端に三次元の共通核を持つ。この端は位置0と隣接し、
$K$の共通核$X$が非零なので、そこにも同じ平面を照合する零点がある。
既に排除した三次元解空間に帰着する。

こうして照合される両束のいずれの端にも二次元以上の共通核はない。
実際、もう一方の束の共通核が非零な端はその両端と隣接する。
簡約後の両端の次元は従って少なくとも3であり、列挙した型でこれを満たすのは$(3)$だけである。
一点のファイバーは射影線型空間の積なので被約でもある。
\end{proof}

\subsection{動く平面の曲面と第三階数4の排除}
\begin{lemma}[動く照合直線が掃く曲面]\label{lem:matched-ruling-quadric}
上の階数4のAS仮定の下で、辺0の照合平面の動く既約族が
$\mathbb P(H)$に掃く曲面は滑らかな二次曲面である。
その直線族は一方の母線族の全体であり、直線ごとに位置0で固定される核点は
三次の有理正規曲線を描く。
\end{lemma}
\begin{proof}
係数行列を$A(z):V_0^*\to V_2$と読む。各照合直線では全点で階数3、
位置0の核は固定であり、余核の双対直線は$L_3^{\mathsf T}$の三次小行列式により
$\mathbb P(V_2^*)$の有理正規三次曲線を動く。
従って$\det A$は恒等的に零ではない。実際、固定核を$x$、余核の双対を$y(z)$とすると、
行列式の任意の方向$z'$への微分は非零倍を除いて$y(z)A(z')x$である。
全て零なら、$y(z)$が$V_2^*$を張るので全ての$A(z')$が$x$を消し、
一因子階数に反する。

同じ向きの異なる照合直線は交わらない。交点での階数3行列の核$x$は一意で、
$x$の評価の核平面も一意だからである。射影的な正次元の既約族から
既約射影曲線を取り、その正規化を$T$とする。ここではGrassmann多様体内の
曲線を取る。固定核$x$は$T$上の射である（束上の共通核は定数階数）。
その平面束の射影化$\pi:Y\to T$を$\mathbb P(H)$へ写すと、像は
$\det A=0$に入る既約曲面$D$である。相異なる直線の非交叉と標数零から
$Y\to D$は双有理であり、超平面の引戻しを$h$とすると$\int_Y h^2=\deg D$である。
同じ非交叉性から照合直線の既約族の次元自体も高々1である。

$d=\pi^*c_1(x^*\mathcal O(1))$とする。$Y$上で
$4\mathcal O(-h)\to4\mathcal O$の核は第一Chern類$-h-d$の直線束、
余核も直線束でその第一Chern類は$3h-d$である。$d^2=0$を使い、
核と余核の完全列からChern指標の次数2部分を比較すると
\[
 (3h-d)^2=-4h^2+(h+d)^2,\qquad 3\deg D=2\deg_T x^*\mathcal O(1).
\]
従って$\deg D$は2または4である。

次数4なら行列式曲面は既約でその台は$D$である。
位置2でも正次元の照合族があるので、その直線族も同じ$D$を掃く。
こちらは余核側が固定され、$L_3$型により核点は動くから、二族に共通の直線はない。
第二族の任意の直線は第一族の全ての直線に会う。
実際、この直線は$D$の中にあり、その各点は第一族の直線上にある。
ただ一本の第一族の直線には収まらず、会う直線のパラメータ集合は閉なので、曲線全体となる。
第一族の相異なる三直線は互いにねじれの位置にあり、これら全てに会う直線は
滑らかな二次曲面上に入る。これは最初の二直線を相補二平面の射影、第三をその間の
同型のグラフとして読むと得られる。従って第二族の掃く$D$は二次曲面に入り、次数4に反する。
ゆえに$D$は既約二次曲面である。円錐では全直線が頂点を通るため、$D$は滑らかである。
射影的な動く族はその一方の母線族$\mathbb P^1$全体を与える。

$D=\mathbb P^1_p\times\mathbb P^1_q$で第一変数$p$を固定核のパラメータとする。
上の等式から$x^*\mathcal O(1)$の次数は3であり、余核直線束は
$\mathcal O(3,3)\otimes\mathcal O(-3,0)=\mathcal O(0,3)$である。
従って余核点は$q$だけで決まる。$q$を固定する直線上では余核が固定で常に階数3なので、
共通零行を除いた$3\times4$束は$L_3$である。その核点$x(p)$は有理正規三次曲線を描く。
同様に$q$側の余核点も有理正規三次曲線である。
\end{proof}

\begin{theorem}[標準AS対象の第三階数は5以上]\label{thm:as-opposite-rank-five-lower}
標準14641型の全AS対象について$5\le\rho\le16$である。
\end{theorem}
\begin{proof}
第三階数4を仮定する。辺0の正次元の照合から前補題の二次曲面を取り、
$H=P\otimes Q$と書いて分解可能元$z=p\otimes q$のSegre像をその曲面とする。
$A(p,q)$は両変数に線型で、全ての点で階数3である。核点は$p$だけ、
余核点は$q$だけで決まる。双対係数空間$H^*=P^*\otimes Q^*$の対応する直線は
$(kp)^\perp\otimes Q^*$なので、こちらも双対Segre曲面の一方の母線族の全体を掃く。
$B(p',q'):V_1^*\to V_3$もそこで全点階数3、核点は$p'$のみに依存する。
この曲面上でも、先の核・余核の直線束の計算と母線への制限により
二つの点は各変数の有理正規三次曲線を描く。

この条件を満たす$A$の表示を明示する。標準三次曲線を核とする
$3\times4$の線型束$D(p)$と、相手の三次曲線を余核点とする
$4\times3$の線型束$M(q)$をそれぞれ$L_3,L_3^{\mathsf T}$から取る。
$D(p)$は$4\mathcal O(-1,-1)\twoheadrightarrow3\mathcal O(0,-1)$で
$A$と核を共有し、$M(q)$は$3\mathcal O(0,-1)\hookrightarrow4\mathcal O$で
$A$と像を共有する。従って$A(p,q)=M(q)C D(p)$であり、
中央の自己同型$C$は定数可逆$3\times3$行列である。
同様に$B(p',q')=M'(q')C'D'(p')$である。
これらは分解可能元での表示だが、それが係数空間を張るので全係数が決まる。

$P,Q$の基底と双対基底で係数を読み、位置0,1から2,3への切断行列を取ると
\[
 \sum_{i,j} A_{ij}\otimes B^{ij}
 =\Big(\sum_j M_j\otimes M'{}^{j}\Big)(C\otimes C')
   \Big(\sum_i D_i\otimes D'{}^{i}\Big).
\]
左の因子は九次元から十六次元への単射である。
実際、$L_3^{\mathsf T}$の二係数は二変数二次式に各変数を掛ける二つの写像である。
従ってテンソル積の和は、別々の二組の変数について非零な双一次式による乗法になる。
変数の基底変更もこの積の非零性を保つ。右の因子は転置で同じ議論を適用して全射である。
ゆえに切断の階数は9であり、標準の隣接切断階数6に反する。
定理\ref{thm:opposite-rank-at-least-four}と合わせて主張を得る。
\end{proof}

\subsection{有限な三つの辺と階数5の残余}
非零$x\in V_i^*$に対し、$R_i$を$x$で縮約した像の次元を$f_i(x)$、
$R_{i-1}$の縮約像の次元を$b_i(x)$と書く。この二数は三テンソル$w(x)$の
巡回順での先頭・末端の支持次元である。$\Gamma_i$が有限なら
どの$x$でも$f_i(x)\ge3$であり、$(x,y)\in\Gamma_i$なら
$f_i(x)=b_{i+1}(y)=3$である。零点の片方を固定したファイバーが線型空間であるためである。

\begin{proposition}[連続三辺の有限性と第三階数]\label{prop:three-finite-edges}
標準AS対象の連続する三つの$\Gamma$が有限ならば$\rho\ge6$である。
従って第三階数5では少なくとも二辺の点対応が正次元である。
\end{proposition}
\begin{proof}
$(x,y)\in\Gamma_i$では常に
\[
 b_i(x)+f_{i+1}(y)\le6 .
\]
実際、$w(x)$の末端因子を縮約して得る$R_{i+1}$の$b_i(x)$次元部分空間は、
$w(x,y)=0$なので$y$による縮約で消える。
前後の辺がともに有限なら左辺の両数は3以上であり、従って両方3である。
中央の辺$\Gamma_i$も有限なら、$w(x)$、$w(y)$のそれぞれで両端の支持次元は3,3になる。
このとき補題\ref{lem:middle-binary-pencil}より中央の支持次元は2になれず、1も既に排除した。
従って$r_i(x),r_{i+1}(y)\ge3$であり、像と核の不等式が結論を与える。
\end{proof}

\begin{corollary}[四辺が有限な場合の点列]\label{cor:finite-edge-point-sequences}
標準有限Frobenius候補で四つの$\Gamma_i$が全て有限ならば、それぞれ長さ20であり、
共有する文字への射影は$\Gamma_i$と$\Gamma_{i+1}$を同じ部分スキームへ同型に写す。
従って隣接零点対は両方向に一意に延長し、始点を固定した点加群の点列を与える。
\end{corollary}
\begin{proof}
上記の支持次元の不等式と有限性による等号にはAS性を使っていない。
$\Gamma_i$上で$w(x)$の末端縮約の像は各閉点で三次元なので、
局所的に$R_{i+1}$の階数3の直和因子を含む。これは$y$による縮約で零になり、
残りの生成元は三本で足りる。従って$R_{i+1}(y)$の階数4の小行列式は
スキームとして$\Gamma_i$上で消える。各閉点で階数3なので、次の文字$z$は
局所的な線型方程式の核直線から射として一意に決まり、$(y,z)\in\Gamma_{i+1}$となる。
逆順にも同じ操作があり、共有文字$y$から元の文字を復元する一意性で両射は逆である。
同じ核直線の構成により共有文字への射影は階数3以下の小行列式スキームへの閉埋込みである。
長さは六本の双一次式の完全交叉から20である。各文字列の関係は隣接二文字だけなので、
この射の反復が全点列を与える。
\end{proof}

\begin{proposition}[任意の第三階数での階数2評価のファイバー]\label{prop:rank-two-evaluation-fiber}
標準AS対象の任意の位置で、階数2評価の行空間を係数空間の双対Grassmann多様体へ送る写像は
各平面上に高々一つの点を持つ。特に$\rho=5$でも、階数2の点が動けばその平面は動く。
\end{proposition}
\begin{proof}
同じ平面に対応する共変ベクトルが独立に二つあるなら、その線型包$X$で
評価は固定した二次元の係数だけを使い、全ての非零評価の階数が2である。
これを位置0とする。制限した二行列$C_s,C_t:X\to V_2$は先の表で
$(2)$または$(1,1)$の行に属し、出力支持$E$の次元は少なくとも3である。
相手の二係数$D_s,D_t$は補題\ref{lem:middle-binary-pencil}から同じ表の四列で尽くされる。
切断の階数6と一因子階数4を使う表の議論は$\rho$を使わないので、
$(1)$列だけが残る。従って相手には固定文字$u$と二次元の支持$Z$がある。
逆順も許して固定文字を位置1に取る。位置0を$X^*$へ射影した部分は
\[
  u\otimes\bigl(C_s\otimes z_s+C_t\otimes z_t\bigr),
 \qquad Z=\langle z_s,z_t\rangle\subset V_3
\]
である（表示は因子を並べ替えたもの）。

元の四窓の一次空間を$X^*,V_1/ku,V_2,V_3/Z$へ落とすと、
位置0に接する二辺の関係像が消える。
$R_1$は$ku\otimes E$を含むので残る辺1の像次元は高々3である。
$R_2$は位置0を$X$で、位置1を$u$に非零な一つの汎関数で縮約した
二次元部分空間を含み、これは$V_2\otimes Z$に入る。
その次元は全ての非零$x\in X$の評価階数2から従う。辺2の像次元は従って高々4である。
位置1からの三文字の商は少なくとも$3\cdot4\cdot2-3\cdot2-3\cdot4=6$次元で、
位置0の二次元因子が独立に繰り返せる。長さ$4m$の商の次元は少なくとも$12^m$となり矛盾する。
\end{proof}
